% TOP_PROBLEM: {"id": 3332, "title": "Inverse Galois Problem", "category_id": 17, "category": {"id": 17, "name": "group_theory", "display_name": "Group Theory", "description": "Problems about groups, group actions, representations, and related algebraic structures.", "slug": "group-theory", "order_index": 17, "created_at": "2026-07-31T15:26:25.671Z"}, "status": "open", "problem_number": "OPG-3572", "set_id": 12, "statusQualification": "Finite Galois extension is explicit, clarifying the catalogue phrase finite algebraic extension."}
% FIRST_POSTED: 2026-10-01
% ENTRY_KIND: statement_only
% UnsolvedMath ID 3332; source catalogue code OPG-3572
% !TeX program = lualatex
% Definition and sources only; this file is not an LLM solution attempt.
\documentclass[11pt,a4paper]{article}
\usepackage[margin=25mm]{geometry}
\usepackage{fontspec}
\setmainfont{lmroman10-regular.otf}[BoldFont=lmroman10-bold.otf,
 ItalicFont=lmroman10-italic.otf,BoldItalicFont=lmroman10-bolditalic.otf]
\usepackage{amsmath,amssymb,mathtools}
\usepackage{unicode-math}
\setmathfont{latinmodern-math.otf}
\AtBeginDocument{\let\setminus\smallsetminus}
\usepackage{xurl,hyperref}
\hypersetup{colorlinks=true,urlcolor=blue}
\setcounter{secnumdepth}{0}
\setlength{\parindent}{0pt}
\setlength{\parskip}{5pt}
\setlength{\emergencystretch}{3em}
\begin{document}
\section*{Inverse Galois Problem}\label{3332}
{\small UnsolvedMath ID 3332 · OPG-3572 · Group Theory · OpenGarden}
\subsection{Definitions and mathematical statement}
Let $\mathbb Q$ denote the rational field and let $G$ be an arbitrary
finite group. A field extension $L/\mathbb Q$ is finite if $L$ has
finite dimension as a vector space over $\mathbb Q$. It is Galois
if it is normal and separable: every irreducible polynomial over
$\mathbb Q$ with a root in $L$ splits completely in $L$, and those
roots are distinct. Separability is automatic in characteristic zero.

The Galois group $\operatorname{Gal}(L/\mathbb Q)$ is the group,
under composition, of field automorphisms of $L$ fixing every
rational number. The inverse Galois problem asks whether
\[
 \forall\text{ finite groups }G\ \exists\text{ a finite Galois extension }
 L/\mathbb Q:\quad\operatorname{Gal}(L/\mathbb Q)\cong G.
\]
The isomorphism is an isomorphism of abstract groups; it need not
realize a previously specified permutation action. For a Galois
extension, the degree $[L:\mathbb Q]$ equals $|G|$.

Equivalently, one asks for a polynomial $f\in\mathbb Q[x]$ whose
splitting field, the smallest field containing all its roots, has
Galois group isomorphic to $G$. The field or polynomial may depend
on $G$, and the base field must remain $\mathbb Q$.

The supplied statement says ``finite algebraic extension''. Here
``Galois group'' is interpreted in the standard inverse-problem
sense by explicitly requiring a finite Galois extension, rather
than merely specifying the automorphism group of a possibly
nonnormal extension. No ramification restrictions or effective
algorithm are additional requirements.
\subsection{Short English statement}
Does every finite group occur as the Galois group of a finite Galois extension of the rational numbers?
\subsection{Sources}
\begin{itemize}
\item[S1] ULAM.AI and the UnsolvedMath Contributors, supplied problems.json, record 3332 (OPG-3572), OpenGarden collection. Catalogue identity and source transcription; CC BY 4.0.
\url{https://huggingface.co/datasets/ulamai/UnsolvedMath}.
\item[S2] Open Problem Garden, Inverse Galois Problem. Original problem and discussion; the mathematical formulation below is expanded from the supplied source record.
\url{http://www.openproblemgarden.org/op/inverse_galois_problem}.
\item[S3] A. Schmidt and K. Wingberg, Šafarevič’s theorem on solvable groups as Galois groups (1998), a primary treatment of realizations over Q.
\url{https://arxiv.org/abs/math/9809211}.
\end{itemize}
\end{document}
